\section{Failure and Atomicity}

\subsection{Atomic Visibility of Hosts Replacement}

During the final replacement performed by \code{removeHostEntry()}, an observer
of the active hosts pathname sees one of two complete values:

\[
  \mathsf{visible}(H)\in\{H_{old},\Remove{d}(H_{old})\}.
\]

No observer should see a partially generated $\Remove{d}(H_{old})$.  The
implementation creates the temporary file in the same directory, writes the
complete retained contents, calls \code{fflush()} and \code{fsync()}, restores
metadata, closes the stream, and then calls \code{rename()}.

\begin{lstlisting}[language=C,caption={Core atomic replacement sequence}]
temporary_fd = mkstemp(temporary_path);
destination = fdopen(temporary_fd, "w");

/* Write the complete revised file. */

if (fflush(destination) == EOF ||
    fsync(fileno(destination)) == -1 ||
    fchmod(fileno(destination),
           hosts_status.st_mode & 07777) == -1)
    goto cleanup;

if (rename(temporary_path, hosts_path) == -1)
    goto cleanup;
\end{lstlisting}

If any step before \code{rename()} fails, the active hosts contents remain
$H_{old}$ and cleanup unlinks the incomplete temporary file.

\subsection{Atomic Visibility Is Not Full Crash Durability}

The replacement file is synchronized before rename.  The containing directory
is not synchronized after rename.  Consequently, the implementation provides
atomic pathname visibility but does not claim the strongest possible guarantee
that the new directory entry survives immediate power loss.

This distinction can be stated as:

\[
  \mathsf{atomicVisibility}\not\Longrightarrow
  \mathsf{maximumCrashDurability}.
\]

A directory \code{fsync()} after rename would strengthen the durability
contract.

\subsection{Add Is Not a Multi-File Transaction}

The allow list changes before the hosts file is replaced.  A later failure can
therefore leave

\[
  d\in\SetOf{A}\quad\land\quad H=H_{old}.
\]

Likewise, \path{hosts.bak} can be updated even if temporary-file generation or
rename later fails.  Atomic replacement protects the active hosts pathname; it
does not make the complete \code{add} action all-or-nothing.

\subsection{Prep and Restore Copy in Place}

\code{copyFile()} opens its destination with truncation.  The \code{prep} and
\code{restore} copies are therefore outside the atomic-replacement law.  A
failure during either copy can leave a partially written destination.

\begin{observation}
Separating primitive laws from whole-command behavior reveals three different
levels of assurance:
\begin{enumerate}
  \item exact and idempotent content transformation;
  \item atomic replacement of one pathname; and
  \item transactional update of several files.
\end{enumerate}
The current implementation provides the first two for hosts-file removal, but
does not claim the third.
\end{observation}
